%% bare_jrnl.tex
%% V1.4b
%% 2015/08/26
%% by Michael Shell
%% see http://www.michaelshell.org/
%% for current contact information.
%%
%% This is a skeleton file demonstrating the use of IEEEtran.cls
%% (requires IEEEtran.cls version 1.8b or later) with an IEEE
%% journal paper.
%%
%% Support sites:
%% http://www.michaelshell.org/tex/ieeetran/
%% http://www.ctan.org/pkg/ieeetran
%% and
%% http://www.ieee.org/

%%*************************************************************************
%% Legal Notice:
%% This code is offered as-is without any warranty either expressed or
%% implied; without even the implied warranty of MERCHANTABILITY or
%% FITNESS FOR A PARTICULAR PURPOSE! 
%% User assumes all risk.
%% In no event shall the IEEE or any contributor to this code be liable for
%% any damages or losses, including, but not limited to, incidental,
%% consequential, or any other damages, resulting from the use or misuse
%% of any information contained here.
%%
%% All comments are the opinions of their respective authors and are not
%% necessarily endorsed by the IEEE.
%%
%% This work is distributed under the LaTeX Project Public License (LPPL)
%% ( http://www.latex-project.org/ ) version 1.3, and may be freely used,
%% distributed and modified. A copy of the LPPL, version 1.3, is included
%% in the base LaTeX documentation of all distributions of LaTeX released
%% 2003/12/01 or later.
%% Retain all contribution notices and credits.
%% ** Modified files should be clearly indicated as such, including  **
%% ** renaming them and changing author support contact information. **
%%*************************************************************************


% *** Authors should verify (and, if needed, correct) their LaTeX system  ***
% *** with the testflow diagnostic prior to trusting their LaTeX platform ***
% *** with production work. The IEEE's font choices and paper sizes can   ***
% *** trigger bugs that do not appear when using other class files.       ***                          ***
% The testflow support page is at:
% http://www.michaelshell.org/tex/testflow/



\documentclass[journal]{IEEEtran}
%
% If IEEEtran.cls has not been installed into the LaTeX system files,
% manually specify the path to it like:
% \documentclass[journal]{../sty/IEEEtran}


\usepackage{graphicx}



% Some very useful LaTeX packages include:
% (uncomment the ones you want to load)


% *** MISC UTILITY PACKAGES ***
%
%\usepackage{ifpdf}
% Heiko Oberdiek's ifpdf.sty is very useful if you need conditional
% compilation based on whether the output is pdf or dvi.
% usage:
% \ifpdf
%   % pdf code
% \else
%   % dvi code
% \fi
% The latest version of ifpdf.sty can be obtained from:
% http://www.ctan.org/pkg/ifpdf
% Also, note that IEEEtran.cls V1.7 and later provides a builtin
% \ifCLASSINFOpdf conditional that works the same way.
% When switching from latex to pdflatex and vice-versa, the compiler may
% have to be run twice to clear warning/error messages.






% *** CITATION PACKAGES ***
%
%\usepackage{cite}
% cite.sty was written by Donald Arseneau
% V1.6 and later of IEEEtran pre-defines the format of the cite.sty package
% \cite{} output to follow that of the IEEE. Loading the cite package will
% result in citation numbers being automatically sorted and properly
% "compressed/ranged". e.g., [1], [9], [2], [7], [5], [6] without using
% cite.sty will become [1], [2], [5]--[7], [9] using cite.sty. cite.sty's
% \cite will automatically add leading space, if needed. Use cite.sty's
% noadjust option (cite.sty V3.8 and later) if you want to turn this off
% such as if a citation ever needs to be enclosed in parenthesis.
% cite.sty is already installed on most LaTeX systems. Be sure and use
% version 5.0 (2009-03-20) and later if using hyperref.sty.
% The latest version can be obtained at:
% http://www.ctan.org/pkg/cite
% The documentation is contained in the cite.sty file itself.






% *** GRAPHICS RELATED PACKAGES ***
%
\ifCLASSINFOpdf
  % \usepackage[pdftex]{graphicx}
  % declare the path(s) where your graphic files are
  % \graphicspath{{../pdf/}{../jpeg/}}
  % and their extensions so you won't have to specify these with
  % every instance of \includegraphics
  % \DeclareGraphicsExtensions{.pdf,.jpeg,.png}
\else
  % or other class option (dvipsone, dvipdf, if not using dvips). graphicx
  % will default to the driver specified in the system graphics.cfg if no
  % driver is specified.
  % \usepackage[dvips]{graphicx}
  % declare the path(s) where your graphic files are
  % \graphicspath{{../eps/}}
  % and their extensions so you won't have to specify these with
  % every instance of \includegraphics
  % \DeclareGraphicsExtensions{.eps}
\fi
% graphicx was written by David Carlisle and Sebastian Rahtz. It is
% required if you want graphics, photos, etc. graphicx.sty is already
% installed on most LaTeX systems. The latest version and documentation
% can be obtained at: 
% http://www.ctan.org/pkg/graphicx
% Another good source of documentation is "Using Imported Graphics in
% LaTeX2e" by Keith Reckdahl which can be found at:
% http://www.ctan.org/pkg/epslatex
%
% latex, and pdflatex in dvi mode, support graphics in encapsulated
% postscript (.eps) format. pdflatex in pdf mode supports graphics
% in .pdf, .jpeg, .png and .mps (metapost) formats. Users should ensure
% that all non-photo figures use a vector format (.eps, .pdf, .mps) and
% not a bitmapped formats (.jpeg, .png). The IEEE frowns on bitmapped formats
% which can result in "jaggedy"/blurry rendering of lines and letters as
% well as large increases in file sizes.
%
% You can find documentation about the pdfTeX application at:
% http://www.tug.org/applications/pdftex





% *** MATH PACKAGES ***
%
%\usepackage{amsmath}
% A popular package from the American Mathematical Society that provides
% many useful and powerful commands for dealing with mathematics.
%
% Note that the amsmath package sets \interdisplaylinepenalty to 10000
% thus preventing page breaks from occurring within multiline equations. Use:
%\interdisplaylinepenalty=2500
% after loading amsmath to restore such page breaks as IEEEtran.cls normally
% does. amsmath.sty is already installed on most LaTeX systems. The latest
% version and documentation can be obtained at:
% http://www.ctan.org/pkg/amsmath





% *** SPECIALIZED LIST PACKAGES ***
%
%\usepackage{algorithmic}
% algorithmic.sty was written by Peter Williams and Rogerio Brito.
% This package provides an algorithmic environment fo describing algorithms.
% You can use the algorithmic environment in-text or within a figure
% environment to provide for a floating algorithm. Do NOT use the algorithm
% floating environment provided by algorithm.sty (by the same authors) or
% algorithm2e.sty (by Christophe Fiorio) as the IEEE does not use dedicated
% algorithm float types and packages that provide these will not provide
% correct IEEE style captions. The latest version and documentation of
% algorithmic.sty can be obtained at:
% http://www.ctan.org/pkg/algorithms
% Also of interest may be the (relatively newer and more customizable)
% algorithmicx.sty package by Szasz Janos:
% http://www.ctan.org/pkg/algorithmicx




% *** ALIGNMENT PACKAGES ***
%
%\usepackage{array}
% Frank Mittelbach's and David Carlisle's array.sty patches and improves
% the standard LaTeX2e array and tabular environments to provide better
% appearance and additional user controls. As the default LaTeX2e table
% generation code is lacking to the point of almost being broken with
% respect to the quality of the end results, all users are strongly
% advised to use an enhanced (at the very least that provided by array.sty)
% set of table tools. array.sty is already installed on most systems. The
% latest version and documentation can be obtained at:
% http://www.ctan.org/pkg/array


% IEEEtran contains the IEEEeqnarray family of commands that can be used to
% generate multiline equations as well as matrices, tables, etc., of high
% quality.




% *** SUBFIGURE PACKAGES ***
%\ifCLASSOPTIONcompsoc
%  \usepackage[caption=false,font=normalsize,labelfont=sf,textfont=sf]{subfig}
%\else
%  \usepackage[caption=false,font=footnotesize]{subfig}
%\fi
% subfig.sty, written by Steven Douglas Cochran, is the modern replacement
% for subfigure.sty, the latter of which is no longer maintained and is
% incompatible with some LaTeX packages including fixltx2e. However,
% subfig.sty requires and automatically loads Axel Sommerfeldt's caption.sty
% which will override IEEEtran.cls' handling of captions and this will result
% in non-IEEE style figure/table captions. To prevent this problem, be sure
% and invoke subfig.sty's "caption=false" package option (available since
% subfig.sty version 1.3, 2005/06/28) as this is will preserve IEEEtran.cls
% handling of captions.
% Note that the Computer Society format requires a larger sans serif font
% than the serif footnote size font used in traditional IEEE formatting
% and thus the need to invoke different subfig.sty package options depending
% on whether compsoc mode has been enabled.
%
% The latest version and documentation of subfig.sty can be obtained at:
% http://www.ctan.org/pkg/subfig




% *** FLOAT PACKAGES ***
%
%\usepackage{fixltx2e}
% fixltx2e, the successor to the earlier fix2col.sty, was written by
% Frank Mittelbach and David Carlisle. This package corrects a few problems
% in the LaTeX2e kernel, the most notable of which is that in current
% LaTeX2e releases, the ordering of single and double column floats is not
% guaranteed to be preserved. Thus, an unpatched LaTeX2e can allow a
% single column figure to be placed prior to an earlier double column
% figure.
% Be aware that LaTeX2e kernels dated 2015 and later have fixltx2e.sty's
% corrections already built into the system in which case a warning will
% be issued if an attempt is made to load fixltx2e.sty as it is no longer
% needed.
% The latest version and documentation can be found at:
% http://www.ctan.org/pkg/fixltx2e


%\usepackage{stfloats}
% stfloats.sty was written by Sigitas Tolusis. This package gives LaTeX2e
% the ability to do double column floats at the bottom of the page as well
% as the top. (e.g., "\begin{figure*}[!b]" is not normally possible in
% LaTeX2e). It also provides a command:
%\fnbelowfloat
% to enable the placement of footnotes below bottom floats (the standard
% LaTeX2e kernel puts them above bottom floats). This is an invasive package
% which rewrites many portions of the LaTeX2e float routines. It may not work
% with other packages that modify the LaTeX2e float routines. The latest
% version and documentation can be obtained at:
% http://www.ctan.org/pkg/stfloats
% Do not use the stfloats baselinefloat ability as the IEEE does not allow
% \baselineskip to stretch. Authors submitting work to the IEEE should note
% that the IEEE rarely uses double column equations and that authors should try
% to avoid such use. Do not be tempted to use the cuted.sty or midfloat.sty
% packages (also by Sigitas Tolusis) as the IEEE does not format its papers in
% such ways.
% Do not attempt to use stfloats with fixltx2e as they are incompatible.
% Instead, use Morten Hogholm'a dblfloatfix which combines the features
% of both fixltx2e and stfloats:
%
% \usepackage{dblfloatfix}
% The latest version can be found at:
% http://www.ctan.org/pkg/dblfloatfix




%\ifCLASSOPTIONcaptionsoff
%  \usepackage[nomarkers]{endfloat}
% \let\MYoriglatexcaption\caption
% \renewcommand{\caption}[2][\relax]{\MYoriglatexcaption[#2]{#2}}
%\fi
% endfloat.sty was written by James Darrell McCauley, Jeff Goldberg and 
% Axel Sommerfeldt. This package may be useful when used in conjunction with 
% IEEEtran.cls'  captionsoff option. Some IEEE journals/societies require that
% submissions have lists of figures/tables at the end of the paper and that
% figures/tables without any captions are placed on a page by themselves at
% the end of the document. If needed, the draftcls IEEEtran class option or
% \CLASSINPUTbaselinestretch interface can be used to increase the line
% spacing as well. Be sure and use the nomarkers option of endfloat to
% prevent endfloat from "marking" where the figures would have been placed
% in the text. The two hack lines of code above are a slight modification of
% that suggested by in the endfloat docs (section 8.4.1) to ensure that
% the full captions always appear in the list of figures/tables - even if
% the user used the short optional argument of \caption[]{}.
% IEEE papers do not typically make use of \caption[]'s optional argument,
% so this should not be an issue. A similar trick can be used to disable
% captions of packages such as subfig.sty that lack options to turn off
% the subcaptions:
% For subfig.sty:
% \let\MYorigsubfloat\subfloat
% \renewcommand{\subfloat}[2][\relax]{\MYorigsubfloat[]{#2}}
% However, the above trick will not work if both optional arguments of
% the \subfloat command are used. Furthermore, there needs to be a
% description of each subfigure *somewhere* and endfloat does not add
% subfigure captions to its list of figures. Thus, the best approach is to
% avoid the use of subfigure captions (many IEEE journals avoid them anyway)
% and instead reference/explain all the subfigures within the main caption.
% The latest version of endfloat.sty and its documentation can obtained at:
% http://www.ctan.org/pkg/endfloat
%
% The IEEEtran \ifCLASSOPTIONcaptionsoff conditional can also be used
% later in the document, say, to conditionally put the References on a 
% page by themselves.




% *** PDF, URL AND HYPERLINK PACKAGES ***
%
%\usepackage{url}
% url.sty was written by Donald Arseneau. It provides better support for
% handling and breaking URLs. url.sty is already installed on most LaTeX
% systems. The latest version and documentation can be obtained at:
% http://www.ctan.org/pkg/url
% Basically, \url{my_url_here}.




% *** Do not adjust lengths that control margins, column widths, etc. ***
% *** Do not use packages that alter fonts (such as pslatex).         ***
% There should be no need to do such things with IEEEtran.cls V1.6 and later.
% (Unless specifically asked to do so by the journal or conference you plan
% to submit to, of course. )


% correct bad hyphenation here
\hyphenation{op-tical net-works semi-conduc-tor}


\begin{document}
%
% paper title
% Titles are generally capitalized except for words such as a, an, and, as,
% at, but, by, for, in, nor, of, on, or, the, to and up, which are usually
% not capitalized unless they are the first or last word of the title.
% Linebreaks \\ can be used within to get better formatting as desired.
% Do not put math or special symbols in the title.
\title{Towards Language-Based Modulation of Assistive Robots through Multimodal Models}
%
%
% author names and IEEE memberships
% note positions of commas and nonbreaking spaces ( ~ ) LaTeX will not break
% a structure at a ~ so this keeps an author's name from being broken across
% two lines.
% use \thanks{} to gain access to the first footnote area
% a separate \thanks must be used for each paragraph as LaTeX2e's \thanks
% was not built to handle multiple paragraphs
%

\author{Philipp~Wicke,
        L\"{u}fti~Kerem~\c{S}enel,
        Shengqiang~Zhang,
        Luis~Figueredo,
        Abdeldjallil~Naceri,
        Sami~Haddadin
        and~Hinrich~Sch\"{u}tze% <-this % stops a space
\thanks{P. Wicke, L.K. \c{S}enel, S. Zhang and H. Sch\"{u}tze are with 
the Center for Information and Language Processing (CIS) and afflilated 
with the Munich Center for Machine Learning (MCML) at the Ludwig-Maximilians-
University (LMU) Munich, Germany. E-mail: \texttt{\{pwicke, lksenel,  shengqiang\}@cis.lmu.de}
}% <-this % stops a space
\thanks{L. Figueredo, A. Naceri and S. Haddadin are with 
the Munich Institute of Robotics and Machine Intelligence (MIRMI)
at Technical University Munich (TUM), Germany. E-mail:
\texttt{\{luis.figueredo,djallil.naceri,haddadin\}@tum.de}}% <-this % stops a space
}

% note the % following the last \IEEEmembership and also \thanks - 
% these prevent an unwanted space from occurring between the last author name
% and the end of the author line. i.e., if you had this:
% 
% \author{....lastname \thanks{...} \thanks{...} }
%                     ^------------^------------^----Do not want these spaces!
%
% a space would be appended to the last name and could cause every name on that
% line to be shifted left slightly. This is one of those "LaTeX things". For
% instance, "\textbf{A} \textbf{B}" will typeset as "A B" not "AB". To get
% "AB" then you have to do: "\textbf{A}\textbf{B}"
% \thanks is no different in this regard, so shield the last } of each \thanks
% that ends a line with a % and do not let a space in before the next \thanks.
% Spaces after \IEEEmembership other than the last one are OK (and needed) as
% you are supposed to have spaces between the names. For what it is worth,
% this is a minor point as most people would not even notice if the said evil
% space somehow managed to creep in.



% The paper headers
\markboth{Geriatronics Summit 2023, July 02 - 03,
Garmisch-Partenkirchen Congress Center}%
{Geriatronics Summit 2023}
% The only time the second header will appear is for the odd numbered pages
% after the title page when using the twoside option.
% 
% *** Note that you probably will NOT want to include the author's ***
% *** name in the headers of peer review papers.                   ***
% You can use \ifCLASSOPTIONpeerreview for conditional compilation here if
% you desire.




% If you want to put a publisher's ID mark on the page you can do it like
% this:
%\IEEEpubid{0000--0000/00\$00.00~\copyright~2015 IEEE}
% Remember, if you use this you must call \IEEEpubidadjcol in the second
% column for its text to clear the IEEEpubid mark.



% use for special paper notices
%\IEEEspecialpapernotice{(Invited Paper)}




% make the title area
\maketitle

% As a general rule, do not put math, special symbols or citations
% in the abstract or keywords.
\begin{abstract}
  In the field of Geriatronics, enabling effective and transparent communication between humans 
  and robots is crucial for enhancing the acceptance and performance of assistive robots. 
  Our early-stage research project investigates the potential of language-based modulation as a 
  means to improve human-robot interaction. We propose to explore 
  real-time modulation during task execution, leveraging language cues, visual 
  references, and multimodal inputs. By developing transparent and interpretable 
  methods, we aim to enable robots to adapt and respond to language commands, enhancing 
  their usability and flexibility. Through the exchange of insights and knowledge at 
  the workshop, we seek to gather valuable feedback to advance our research and 
  contribute to the development of interactive robotic systems for Geriatronics 
  and beyond.
\end{abstract}

% Note that keywords are not normally used for peerreview papers.
\begin{IEEEkeywords}
  HRI, Geriatronics, NLP, Task Modulation
\end{IEEEkeywords}






% For peer review papers, you can put extra information on the cover
% page as needed:
% \ifCLASSOPTIONpeerreview
% \begin{center} \bfseries EDICS Category: 3-BBND \end{center}
% \fi
%
% For peerreview papers, this IEEEtran command inserts a page break and
% creates the second title. It will be ignored for other modes.
%\IEEEpeerreviewmaketitle



\section{Introduction}
In the field of Geriatronics, where robots play a crucial role in assisting 
older adults with daily tasks, effective human-robot interaction is paramount. 
Imagine a scenario where an older adult needs assistance with taking medication. 
The robot is instructed to ``bring the pill bottle from the kitchen counter.'' 
However, what if the user realizes they need to clarify their request by adding, 
``not the red, but the white bottle''? This language-based modulation is a fundamental 
aspect of language-mediated robot control that can significantly enhance the usability 
and adaptability of robots in real-world scenarios.

Addressing the limitations of language-mediated robot control and the need for 
modulation, this research project aims to explore and develop methods for effective 
modulation in human-robot interaction. Its first focus is on creating a dataset for supporting 
machine learning research on modulation. To achieve this, a diverse dataset will be constructed 
through crowdsourcing, incorporating a range of household tasks relevant to Geriatronics. 
Instructions for modulations will be collected, encompassing both low-level modifications 
that affect specific parameters during task execution and high-level modifications that 
involve adding, removing, or modifying tasks within a larger context. The dataset will 
include instructions with varying levels of specificity to enable the training of 
sophisticated data-driven models. In a second step, the focus shifts towards the design
of transparent methods that leverage deep learning techniques with multimodal models
while providing explanations for the robot's behavior. 

This early-stage research project presents the foundation 
required to complete the envisioned scenarios of language-based modulation
in the context of Geriatronics and beyond.

\section{Related Works and Background}

We provide a comprehensive overview of the research landscape related to
 language-mediated modulation in the context of human-robot interaction. 
 This section is divided into two parts to address distinct aspects of the topic. 
 The first part focuses on the need for a benchmark dataset specifically designed 
 for language-mediated modulation tasks, highlighting the limitations of 
 existing datasets. The second part introduces a transparent approach that 
 incorporates real-time segmentation to enhance the understanding of robot 
 actions during modulated tasks and the robot's task performance itself.

\paragraph*{Language-Mediated Modulation}
In the field of human-robot communication, the ability to modulate robot actions 
using natural language is crucial. However, there is currently a lack of rich 
datasets that specifically address language-mediated modulation. To address this 
gap, our project focuses on creating a benchmark dataset for language-mediated 
modulation and sharing it with the research community. This dataset will enable 
the development of multimodal machine learning approaches for improving human-robot 
interaction.

Existing datasets in the field of robot learning and language-conditioned policies 
provide valuable resources for training models to perform various household tasks. 
On the one hand, datasets such as LangLFP \cite{lynch2021langlfp}, LOREL \cite{nair2022lorel}, 
BC-Z \cite{jang2022bcz}, SayCan \cite{brohan2022saycan}, RT-1 \cite{brohan2022rt1}, 
language-table \cite{lynch2022language} or Latte \cite{bucker2022latte} consist of 
labeled trajectories with language annotations, enabling the learning of 
language-conditioned policies for robot control in performing various low-level 
tasks, such as ``navigate to the kitchen'', ``place the book on the table'', or ``avoid 
obstacles on the floor''.

On the other hand, datasets such as ALFRED \cite{shridhar2020alfred}, 
DialFRED \cite{gao2022dialfred}, Playhouse \cite{team2021play}, TEACh 
\cite{padmakumar2022teach}, and TEAChDA \cite{gella2022teachda} are oriented 
towards accomplishing complex objectives at a higher level, involving the 
learning of action plans with multiple substeps, such as ``prepare a meal'' or
``assemble furniture''. They encompass multiple modalities, integrating visual 
and language information. The degree of interaction between humans and robots
varies among these datasets, with some primarily focusing on the initial 
instruction without extensive ongoing interaction. However, most of these datasets 
do not consider language-mediated modulation specifically, and their focus is
either on single-step tasks or multi-step action plans without explicit 
interaction or modulation during execution.

\paragraph*{Transparent Approach Using Real-Time Segmentation}

\begin{figure*}[t!]
    \centering
    \includegraphics[width=\textwidth]{fig01.png}
    \caption{Demonstration of the transparent approach
    for language-based modulation. The initial task
    \textit{stack the cups} is modulated with the command
    \textit{stack object 2 first!} These
    object labels are available to the user in real-time 
    throughout the task execution. For the underlying model, 
    segmentation and label of the target object become available upon 
    modulation for all different camera perspectives.
    \label{fig:demo}}
\end{figure*}


The second part of the proposed research focusses on a
transparent approach for language-based modulation that enables 
robust reference to environmental elements in human-robot communication.
In this context, modulation involves modifying previous 
commands or ongoing actions, requiring the robot to establish 
references to facilitate effective interaction with (elderly) individuals. 
To address this, we aim to create a shared scene representation that both 
humans and robots can access and interpret, facilitating accurate 
references and enhancing communication and interaction for both sides.

Related works have explored the concept of transparency in 
language-mediated robot control. For instance, researchers 
have demonstrated the utility of object labeling in facilitating 
explicit learning or differentiation of novel objects 
\cite{hirschmanner2021investigating,schuetz2017spoken,tucker2020learning}. 
In the context of Geriatronics, this could involve labeling and 
recognizing various everyday objects used by elderly individuals, 
such as medication bottles, walkers, or assistive devices. 
Additionally, transparent methods have been employed to enhance 
communication between robots and humans in various settings 
\cite{bucker2022reshaping,liu2017ar,wicke2021are}. 
For example, a spatial referencing system for mobile robots allows users 
to add waypoints using augmented reality, enabling precise navigation 
assistance \cite{chacko2020augmented}. 
Similarly, in an industrial robot context, visual tracking and 
automated labeling of objects ensure transparent communication through shared labels \cite{liu2017ar}.

Furthermore, human-robot dialogue has been explored as a means to clarify 
ambiguous references and improve language understanding 
\cite{vogel2010eye,parde2015grounding,thomason2016learning,yang2018visual}. 
These studies focus on learning to ground object attributes and names through 
dialogue interactions, which could be applicable to Geriatronics scenarios 
where the robot needs to understand and interpret specific instructions 
related to elderly individuals' needs and preferences. Additionally, 
interactive systems that allow robots to ask follow-up clarifications to 
disambiguate described objects have shown improvements in correctly identifying 
objects with fewer attempts \cite{dogan2022asking}. This approach could be beneficial in situations 
where a person provides instructions to a robot and the robot seeks 
clarification to ensure accurate actions.

Our proposed approach to transparency in language-based modulation differs 
from previous works in two core aspects. Firstly, the language-based interaction 
directly modulates the robot's ongoing actions, aiming to provide immediate 
assistance to individuals. Secondly, our shared scene 
representation serves dual purposes: updating humans about the robot's 
understanding of the environment and supporting natural and robust modulation 
commands. This allows references in natural language commands, such as 
specifying a particular medication or indicating the location of an object,
to be directly linked to objects in the shared scene representation. 
Given that misunderstandings may be critical, this transparency can improve
the user's trust in the system and the interaction. 

\section{Proposed Framework}

\paragraph*{Robotic Environment}


We will utilize RLBench \cite{james2020rlbench} to build our 
dataset by curating and augmenting tasks from its task library.
RLBench offers a convenient platform for dataset creation due to 
its accessibility and task customization capabilities. 
Specifically, it provides automatically generated demonstrations
for a diverse set of one hundred household tasks, which we can 
augment with modulations for our dataset.

RLBench supports multiple variations of tasks, allowing us to 
introduce variability in factors such as target objects and 
colors. Moreover, it enables the generation of unlimited 
episodes for each task variation, enabling us to vary object 
positions across episodes. By incorporating high and low-level 
modulations into the RLBench tasks, we can further diversify 
the dataset. Modulation points can be strategically introduced 
at different stages of task execution, capturing a range 
of modulation scenarios, adding further complexity to the
dataset.

Using RLBench as our dataset generation environment offers 
several advantages. Firstly, RLBench provides a comprehensive 
task library that covers various household tasks, making it 
suitable for studying language-based modulation in the context 
of human-robot interaction. Secondly, the availability of 
automatically generated demonstrations saves time and effort 
in dataset creation. Lastly, the flexibility of RLBench allows 
for easy customization of tasks and the introduction of 
modulations, facilitating the creation of a diverse 
dataset.

\subsection{Dataset Collection and Curation}

Our dataset will incorporate paired demonstrations of low- 
and high-level modulations accompanied by natural language 
instructions, with a possible focus on applications in Geriatronics. 
To gather these instructions, we will employ crowdsourcing, aiming 
to collect approximately 50 natural language instructions for each 
modulation in 30 different household-related tasks. Four conditions
will be considered, resulting from two levels of modulation 
(low- and high-level) and two specificity levels of instructions
(low- and high-specificity).

Low-level (LL) modulations involve modifying a specific parameter 
of an action as it occurs in the Geriatronics context. For instance, 
the instruction ``hand me the pill bottle'' can be LL-modulated by 
saying ``not the brown, but the white one'' or ``be gentle to move it.''
On the other hand, high-level (HL) modulations involve specifying, 
modifying, adding, or removing low-level tasks within the context 
of a high-level task. HL modulations serve to address vagueness in 
the wording of the modulated command or to modify the command itself. 
For example, the command ``bring the cane'' may require the 
modulation command ``avoid stepping on the carpet'' when the robot 
starts walking across the room.

Two types of instructions will be collected: high-specificity 
(HS) and low-specificity (LS). HS modulations provide a detailed 
description of the desired effect of the modulation, leaving 
little ambiguity. Conversely, LS modulations are brief, vague, 
or ambiguous modulations, relying mostly on the information 
provided by the high-level task instruction. 

Crowdsourcing involves presenting crowdworkers with two videos: 
one depicting the robot performing the task without modulation 
and another showing the task with a visual cue indicating modulation.
The crowdworkers are then asked to provide two commands that 
could trigger the modulation: a brief and natural command (LS) 
and a precise and detailed command (HS). Importantly, 
crowdworkers only observe the effect of modulation on 
the robot's behavior and do not receive a modulation 
command directly, as they provide the command.

To facilitate crowdsourcing, the videos include a red exclamation 
mark as a ``modulation point,'' indicating the moment where the 
robot's behavior deviates from its original execution. 
Multiple exclamation marks may be used for tasks with multiple 
modulations. Crowdworkers are requested to rate the difficulty of 
generating the appropriate command and provide an explanation if 
they were unable to generate a command.

For each modulation instance corresponding to a unique modulation 
video in the Geriatronics domain, we will collect a minimum of 30 
high-specificity and low-specificity modulation commands. Leveraging
the capabilities of RLBench tasks, we can create multiple modulation
instances, ensuring the generation of many, diverse data points.

\subsection{Transparent Approach for Language-Based Modulation}

Automatic computation of accurate and semantically rich object
labels is one of the long-term goals of vision research. 
While much progress has been made, the state of the art of 
this technology (e.g., YOLO \cite{wang2022yolo}, DETR 
\cite{carion2020detr}) is not always robust. 
More importantly, even if a semantic label is correct 
(e.g., ``book,'' ``plate''), in many cases, there are multiple
objects present with the same broad semantic label, making
it difficult for robots to identify the correct object. 
For example, it can be challenging for a system to identify 
``medication'' as a label. While human-human communication 
relies on identifying salient differentiating attributes 
(``the blue pill bottle,'' ``the prescription box''), reference to 
objects with differentiating attributes when communicating with 
humans and robots is a field of ongoing, fast-developing research.

Our objective is therefore not to mimic how humans refer to 
objects. We do not want to produce the type of ``natural'' labels 
that humans would use for communication. Instead, we will attempt 
to create simple and distinguishing labels that can be easily 
shared by humans and robots and that can be easily linked across 
language and visual modalities. As an example, consider a scenario 
where a user provides a vague command such as ``Fetch the 
medication from the shelf.'' While our approach may not be 
able to label the object as ``medication,'' the user can refer 
to a label provided by a real-time interface of the robot's 
view, ensuring that the robot retrieves the correct target 
object from the shelf. This is a clear advantage of our 
system, as it allows for modulation of the command and 
improves the accuracy of the robot's response in assisting 
individuals with their medication management.

The proposed models for addressing the RLBench benchmark 
consist of two main components: an input encoder that handles 
visual input from cameras and textual input from user instructions 
to produce encoded representations, and a transformer policy that
utilizes the encoded representations to predict the next robot 
action, typically described in terms of the robot's gripper's 
position, quaternion, and state. Our research builds upon these 
architectures by introducing an augmentation module to enhance 
textual and visual inputs. This augmentation aims to establish 
a clear reference between the object mentioned in the human 
command and a target object segmented or highlighted in the
robot's visual input.

To establish this reference, two key components are employed: 
a shared scene representation and unique synthetic labels
(see Fig. \ref{fig:demo} for an example). 
By using a shared scene representation, we provide a consistent 
understanding of the environment. Synthetic labels are introduced 
to replace potentially ambiguous references, allowing users to 
use specific identifiers such as ``object \#1'' to refer to target 
objects. To facilitate this, users are presented with an augmented 
real-time video feed of the robot's visual inputs, where each 
detected object is labeled with a synthetic label. These synthetic 
labels serve as a clear reference for an object named in the 
command.

When a user refers to a target object using its synthetic label, 
the augmentation module modifies the textual and visual inputs. 
The synthetic label in the textual input is replaced with a special 
token ``TARGET'', while the corresponding object in the visual input
is highlighted or marked. The specific method of marking the referred
object in the visual input, such as using bounding boxes or 
object segmentation boundaries, will be determined based on 
state-of-the-art visual object detection models. During training, 
we hypothesize that the multimodal encoder will learn to associate 
the special token ``TARGET'' with the marked object in the visual 
input, enabling the policy transformer to utilize this association 
for action prediction.

To ensure consistent object tracking across different camera angles 
and time steps, related work and state-of-the-art models in object 
detection and tracking are leveraged. We also consider the 
possibility of assigning unique synthetic labels to objects 
across different camera angles, allowing the augmentation of 
visual input for a specific camera angle containing the mentioned 
synthetic label. This exploration is essential as communication 
robustness improves when a single camera angle is sufficient for 
resolving object references.

In terms of data collection, we will adapt the crowdsourced 
dataset to include synthetic labels. This eliminates the need
for separately collecting instructions with synthetic 
labels. 

\section{Current Progress}

This section presents an overview of the current progress of our 
research, focusing on the successful setup of the RLBench environment
and the implementation of RLBench household tasks. Through the utilization
of this environment, we have identified specific tasks that lend themselves
well to the incorporation of language-based modulations, opening avenues 
for investigating the efficacy of our proposed approach.

To validate the feasibility of our approach, we initially concentrate our 
efforts on a single task as a proof of concept. Specifically, we selected 
the task ``stack the cups'' as our initial testbed. This task involves the
robot stacking two cups into a third cup, thereby establishing a tangible 
scenario for evaluating the effectiveness of language-based modulations in
enhancing robotic performance. Although we have identified many more complex
tasks that allow for modulation, we select a task with relatively short duration
and complexity to reduce computing and training times in this early
phase of the project.

In our initial experimentation, we introduced a modulation instruction in the 
form of ``stack the other cup first.'' This modulation aims to guide the robot's
actions by explicitly specifying the order in which the cups should be stacked. 
By incorporating language-based modulations into the task, our objective is to 
improve the robot's understanding of the task and enhance its ability
to reliably execute complex actions.

To assess the performance of our approach, we train a small-scale model with the aim of
solving the ``stack the cups'' task. The training process is in-line with models by \cite{james2020rlbench}. 
By focusing on a manageable task scale, we seek to evaluate the feasibility and 
effectiveness of our approach before scaling up to more intricate tasks, which
require vastly more computation.

Our current model remains relatively limited in terms of its capacity and capabilities, but
nonetheless allows us to collect images for all camera angles in order to implement the 
transparent approach. Moving forward, our immediate objective is to incorporate textual 
and image augmentation techniques to examine whether the proposed transparent approach further 
enhances the performance and reliablity of our model. 

At this critical juncture, feedback from the Geriatronics community becomes invaluable. 
We recognize the significance of soliciting input and insights from domain experts and
practitioners to refine our approach. The feedback obtained will guide us in iteratively
refining our methodology, ensuring that our proposed research aligns with the specific requirements
and challenges faced in the field of Geriatronics.

\section{Conclusion}

In conclusion, our research in language-based modulations for robotic systems 
demonstrated the potential of natural language instructions to enhance robot adaptability
and responsiveness. Through the use of the RLBench environment, we plan to validate the feasibility 
of integrating language-based modulations into robotic tasks, as exemplified by our proof-of-concept on 
the ``stacking cups'' task (Figure \ref{fig:demo}).

Collaboration with the Geriatronics community is vital in refining our approach and
addressing specific challenges in assistive robotics. With further research and 
development, we envision a future where assistive robots seamlessly integrate natural
language instructions, providing tailored support and enhancing the quality of life
for the elderly in Geriatronics applications.


% An example of a floating figure using the graphicx package.
% Note that \label must occur AFTER (or within) \caption.
% For figures, \caption should occur after the \includegraphics.
% Note that IEEEtran v1.7 and later has special internal code that
% is designed to preserve the operation of \label within \caption
% even when the captionsoff option is in effect. However, because
% of issues like this, it may be the safest practice to put all your
% \label just after \caption rather than within \caption{}.
%
% Reminder: the "draftcls" or "draftclsnofoot", not "draft", class
% option should be used if it is desired that the figures are to be
% displayed while in draft mode.
%
%\begin{figure}[!t]
%\centering
%\includegraphics[width=2.5in]{myfigure}
% where an .eps filename suffix will be assumed under latex, 
% and a .pdf suffix will be assumed for pdflatex; or what has been declared
% via \DeclareGraphicsExtensions.
%\caption{Simulation results for the network.}
%\label{fig_sim}
%\end{figure}

% Note that the IEEE typically puts floats only at the top, even when this
% results in a large percentage of a column being occupied by floats.


% An example of a double column floating figure using two subfigures.
% (The subfig.sty package must be loaded for this to work.)
% The subfigure \label commands are set within each subfloat command,
% and the \label for the overall figure must come after \caption.
% \hfil is used as a separator to get equal spacing.
% Watch out that the combined width of all the subfigures on a 
% line do not exceed the text width or a line break will occur.
%
%\begin{figure*}[!t]
%\centering
%\subfloat[Case I]{\includegraphics[width=2.5in]{box}%
%\label{fig_first_case}}
%\hfil
%\subfloat[Case II]{\includegraphics[width=2.5in]{box}%
%\label{fig_second_case}}
%\caption{Simulation results for the network.}
%\label{fig_sim}
%\end{figure*}
%
% Note that often IEEE papers with subfigures do not employ subfigure
% captions (using the optional argument to \subfloat[]), but instead will
% reference/describe all of them (a), (b), etc., within the main caption.
% Be aware that for subfig.sty to generate the (a), (b), etc., subfigure
% labels, the optional argument to \subfloat must be present. If a
% subcaption is not desired, just leave its contents blank,
% e.g., \subfloat[].


% An example of a floating table. Note that, for IEEE style tables, the
% \caption command should come BEFORE the table and, given that table
% captions serve much like titles, are usually capitalized except for words
% such as a, an, and, as, at, but, by, for, in, nor, of, on, or, the, to
% and up, which are usually not capitalized unless they are the first or
% last word of the caption. Table text will default to \footnotesize as
% the IEEE normally uses this smaller font for tables.
% The \label must come after \caption as always.
%
%\begin{table}[!t]
%% increase table row spacing, adjust to taste
%\renewcommand{\arraystretch}{1.3}
% if using array.sty, it might be a good idea to tweak the value of
% \extrarowheight as needed to properly center the text within the cells
%\caption{An Example of a Table}
%\label{table_example}
%\centering
%% Some packages, such as MDW tools, offer better commands for making tables
%% than the plain LaTeX2e tabular which is used here.
%\begin{tabular}{|c||c|}
%\hline
%One & Two\\
%\hline
%Three & Four\\
%\hline
%\end{tabular}
%\end{table}


% Note that the IEEE does not put floats in the very first column
% - or typically anywhere on the first page for that matter. Also,
% in-text middle ("here") positioning is typically not used, but it
% is allowed and encouraged for Computer Society conferences (but
% not Computer Society journals). Most IEEE journals/conferences use
% top floats exclusively. 
% Note that, LaTeX2e, unlike IEEE journals/conferences, places
% footnotes above bottom floats. This can be corrected via the
% \fnbelowfloat command of the stfloats package.




% Can use something like this to put references on a page
% by themselves when using endfloat and the captionsoff option.
\ifCLASSOPTIONcaptionsoff
  \newpage
\fi



% trigger a \newpage just before the given reference
% number - used to balance the columns on the last page
% adjust value as needed - may need to be readjusted if
% the document is modified later
%\IEEEtriggeratref{8}
% The "triggered" command can be changed if desired:
%\IEEEtriggercmd{\enlargethispage{-5in}}

% references section

% can use a bibliography generated by BibTeX as a .bbl file
% BibTeX documentation can be easily obtained at:
% http://mirror.ctan.org/biblio/bibtex/contrib/doc/
% The IEEEtran BibTeX style support page is at:
% http://www.michaelshell.org/tex/ieeetran/bibtex/
%\bibliographystyle{IEEEtran}
% argument is your BibTeX string definitions and bibliography database(s)
%\bibliography{IEEEabrv,../bib/paper}
%
% <OR> manually copy in the resultant .bbl file
% set second argument of \begin to the number of references
% (used to reserve space for the reference number labels box)
\begin{thebibliography}{1}

  \bibitem{lynch2021langlfp} C. Lynch and P. Sermanet. Language conditioned imitation learning over unstructured data. Robotics: Science and
  Systems, 2021.
  \bibitem{nair2022lorel} S. Nair, E. Mitchell, K. Chen, S. Savarese, C. Finn, et al. Learning language-conditioned robot behavior from offline
  data and crowd-sourced annotation. In Conference on Robot Learning, pages 1303–1315. PMLR, 2022.
  \bibitem{jang2022bcz} E. Jang, A. Irpan, M. Khansari, D. Kappler, F. Ebert, C. Lynch, S. Levine, and C. Finn. Bc-z: Zero-shot task generalization with robotic imitation learning. In Conference on Robot Learning, pages 991–1002. PMLR, 2022.
  \bibitem{brohan2022saycan} A. Brohan, Y. Chebotar, C. Finn, K. Hausman, A. Herzog, D. Ho, J. Ibarz, A. Irpan, E. Jang, R. Julian, et al. Do as i
  can, not as i say: Grounding language in robotic affordances. In 6th Annual Conference on Robot Learning, 2022.
  \bibitem{brohan2022rt1} A. Brohan, N. Brown, J. Carbajal, Y. Chebotar, J. Dabis, C. Finn, K. Gopalakrishnan, K. Hausman, A. Herzog, J. Hsu,
  et al. Rt-1: Robotics transformer for real-world control at scale. arXiv preprint arXiv:2212.06817, 2022.
  \bibitem{lynch2022language} C. Lynch, A. Wahid, J. Tompson, T. Ding, J. Betker, R. Baruch, T. Armstrong, and P. Florence. Interactive language:
  Talking to robots in real time. arXiv preprint arXiv:2210.06407, 2022.
  \bibitem{bucker2022latte} A. F. C. Bucker, L. Figueredo, S. Haddadin, A. Kapoor, S. Ma, S. Vemprala, and R. Bonatti. Latte: Language
  trajectory transformer. In Foundation Models for Decision Making @NEURIPS2022, 2022.
  \bibitem{shridhar2020alfred} M. Shridhar, J. Thomason, D. Gordon, Y. Bisk, W. Han, R. Mottaghi, L. Zettlemoyer, and D. Fox. Alfred: A benchmark
  for interpreting grounded instructions for everyday tasks. In Proceedings of the IEEE/CVF conference on computer
  vision and pattern recognition, pages 10740–10749, 2020.
  \bibitem{gao2022dialfred} X. Gao, Q. Gao, R. Gong, K. Lin, G. Thattai, and G. S. Sukhatme. Dialfred: Dialogue-enabled agents for embodied
  instruction following. IEEE Robotics and Automation Letters, 7(4):10049–10056, 2022.
  \bibitem{team2021play} D. I. A. Team, J. Abramson, A. Ahuja, A. Brussee, F. Carnevale, M. Cassin, F. Fischer, P. Georgiev, A. Goldin,
  T. Harley, et al. Creating multimodal interactive agents with imitation and self-supervised learning. arXiv preprint
  arXiv:2112.03763, 2021.
  \bibitem{padmakumar2022teach} A. Padmakumar, J. Thomason, A. Shrivastava, P. Lange, A. Narayan-Chen, S. Gella, R. Piramuthu, G. Tur, and
  D. Hakkani-Tur. Teach: Task-driven embodied agents that chat. In Proceedings of the AAAI Conference on Artifcial Intelligence, volume 36, pages 2017–2025, 2022.
  \bibitem{gella2022teachda} S. Gella, A. Padmakumar, P. Lange, and D. Hakkani-Tur. Dialog acts for task driven embodied agents. In Proceedings
  of the 23rd Annual Meeting of the Special Interest Group on Discourse and Dialogue, pages 111–123, Edinburgh, UK,
  Sept. 2022. Association for Computational Linguistics. 2022.
  \bibitem{hirschmanner2021investigating} M. Hirschmanner, S. Gross, S. Zafari, B. Krenn, F. Neubarth, and M. Vincze. Investigating transparency methods in a
  robot word-learning system and their effects on human teaching behaviors. In 2021 30th IEEE International Conference
  on Robot \& Human Interactive Communication (RO-MAN), pages 175–182. IEEE, 2021.
  \bibitem{schuetz2017spoken} M. Scheutz, E. Krause, B. Oosterveld, T. Frasca, and R. Platt. Spoken instruction-based one-shot object and action learning in a cognitive robotic architecture. In Proceedings of the 16th Conference on Autonomous Agents and MultiAgent
  Systems, 2017.
  \bibitem{tucker2020learning} M. Tucker, D. Aksaray, R. Paul, G. J. Stein, and N. Roy. Learning unknown groundings for natural language interaction
  with mobile robots. In Robotics Research: The 18th International Symposium ISRR, pages 317–333. Springer, 2020.
  \bibitem{bucker2022reshaping} A. Bucker, L. Figueredo, S. Haddadin, A. Kapoor, S. Ma, and R. Bonatti. Reshaping robot trajectories using
  natural language commands: A study of multi-modal data alignment using transformers. In 2022 IEEE/RSJ
  International Conference on Intelligent Robots and Systems (IROS), pages 978–984. IEEE, 2022.
  \bibitem{liu2017ar} H. Liu and L. Wang. An ar-based worker support system for human-robot collaboration. Procedia Manufacturing,
  11:22–30, 2017.
  \bibitem{wicke2021are} P. Wicke and T. Veale. Are you not entertained? computational storytelling with non-verbal interaction. In
  Companion of the 2021 ACM/IEEE International Conference on Human-Robot Interaction, pages 200–204, 2021.
  \bibitem{chacko2020augmented} S. M. Chacko, A. Granado, A. RajKumar, and V. Kapila. An augmented reality spatial referencing system for mobile
  robots. In 2020 IEEE/RSJ International Conference on Intelligent Robots and Systems (IROS), pages 4446–4452. IEEE,
  2020.
  \bibitem{vogel2010eye} A. Vogel, K. Raghunathan, and D. Jurafsky. Eye spy: Improving vision through dialog. In 2010 AAAI Fall Symposium
  Series. Citeseer, 2010.
  \bibitem{parde2015grounding} N. Parde, A. Hair, M. Papakostas, K. Tsiakas, M. Dagioglou, V. Karkaletsis, and R. D. Nielsen. Grounding the meaning of words through vision and interactive gameplay. In Twenty-Fourth International Joint Conference on Artifcial
  Intelligence, 2015.
  \bibitem{thomason2016learning} J. Thomason, J. Sinapov, M. Svetlik, P. Stone, and R. J. Mooney. Learning multi-modal grounded linguistic semantics
  by playing" i spy". In IJCAI, pages 3477–3483, 2016.
  \bibitem{yang2018visual} J. Yang, J. Lu, S. Lee, D. Batra, and D. Parikh. Visual curiosity: Learning to ask questions to learn visual recognition.
  In A. Billard, A. Dragan, J. Peters, and J. Morimoto, editors, Proceedings of The 2nd Conference on Robot Learning,
  volume 87 of Proceedings of Machine Learning Research, pages 63–80. PMLR, 2018.
  \bibitem{dogan2022asking} F. I. Do\u{g}an, I. Torre, and I. Leite. Asking follow-up clarifcations to resolve ambiguities in human-robot conversation. ˘
  In 2022 17th ACM/IEEE International Conference on Human-Robot Interaction (HRI), pages 461–469. IEEE, 2022.
  \bibitem{james2020rlbench} S. James, Z. Ma, D. R. Arrojo, and A. J. Davison. Rlbench: The robot learning benchmark \& learning environment.
  IEEE Robotics and Automation Letters, 5(2):3019–3026, 2020.
  \bibitem{wang2022yolo} C.-Y. Wang, A. Bochkovskiy, and H.-Y. M. Liao. Yolov7: Trainable bag-of-freebies sets new state-of-the-art for realtime object detectors. arXiv preprint arXiv:2207.02696, 2022.
  \bibitem{carion2020detr} N. Carion, F. Massa, G. Synnaeve, N. Usunier, A. Kirillov, and S. Zagoruyko. End-to-end object detection with transformers. In European conference on computer vision, pages 213–229. Springer, 2020.
  
\end{thebibliography}

% biography section
% 
% If you have an EPS/PDF photo (graphicx package needed) extra braces are
% needed around the contents of the optional argument to biography to prevent
% the LaTeX parser from getting confused when it sees the complicated
% \includegraphics command within an optional argument. (You could create
% your own custom macro containing the \includegraphics command to make things
% simpler here.)
%\begin{IEEEbiography}[{\includegraphics[width=1in,height=1.25in,clip,keepaspectratio]{mshell}}]{Michael Shell}
% or if you just want to reserve a space for a photo:



% You can push biographies down or up by placing
% a \vfill before or after them. The appropriate
% use of \vfill depends on what kind of text is
% on the last page and whether or not the columns
% are being equalized.

%\vfill

% Can be used to pull up biographies so that the bottom of the last one
% is flush with the other column.
%\enlargethispage{-5in}



% that's all folks
\end{document}


